% Chapter 4, Section 2 _Linear Algebra_ Jim Hefferon
%  http://joshua.smcvt.edu/linearalgebra
%  2001-Jun-12
\section{Geometri Determinan}
Bagian sebelumnya mengembangkan determinan secara aljabar, dengan
meninjau rumus-rumus yang memenuhi syarat-syarat tertentu.
Bagian ini melengkapinya dengan pendekatan geometris.
Selain daya tarik intuitifnya, pendekatan ini memiliki keunggulan: sementara
sejauh ini kita hanya meninjau apakah suatu determinan bernilai nol atau tidak,
di sini kita akan memberi makna pada nilai determinan tersebut.
(Bagian sebelumnya memperlakukan determinan sebagai fungsi dari
baris-baris, tetapi bagian ini berfokus pada kolom-kolom.)










\subsection{Determinan sebagai Fungsi Ukuran}
Gambar jajaran genjang ini
sudah dikenal dari konstruksi jumlah dua vektor.
\begin{center}
  \includegraphics{det/mp/ch4.30}
\end{center}

\begin{definition} \label{df:Box}
%<*df:Box>
Di $\Re^n$,
\definend{kotak}\index{kotak}
(atau \definend{paralelepiped}\index{paralelepiped})
yang dibentuk oleh
\( \sequence{\vec{v}_1,\dots,\vec{v}_n} \) 
adalah himpunan
\( \set{t_1\vec{v}_1+\dots+t_n\vec{v}_n
      \suchthat t_1,\ldots,t_n\in \closedinterval{0}{1}} \).
%</df:Box>
\end{definition}

\noindent Jadi, jajaran genjang di atas adalah kotak yang dibentuk oleh
$\sequence{\binom{x_1}{y_1},\binom{x_2}{y_2}}$.
Kotak dalam ruang tiga dimensi ditampilkan dalam \nearbyexample{ex:VolParPiped}.

Kita dapat mencari luas kotak di atas dengan menggambar persegi panjang yang
melingkupinya lalu mengurangkan luas daerah-daerah yang tidak termasuk kotak.
\begin{center}
  \parbox{1.5in}{\hbox{}\hfil\includegraphics{det/mp/ch4.31}\hfil\hbox{}}
  \quad
  \parbox{3.0in}{
    \hbox{}\hfil
    $\begin{array}{l}
       \text{luas jajaran genjang}                    \\
         \hbox{}\quad \hbox{}
            =\text{luas persegi panjang}
                -\text{luas $A$}-\text{luas $B$} \\
         \hbox{}\qquad \hbox{}
            -\cdots-\text{luas $F$}                   \\
         \hbox{}\quad \hbox{}
            =(x_1+x_2)(y_1+y_2)-x_2y_1-x_1y_1/2        \\
         \hbox{}\qquad \hbox{}
            -x_2y_2/2-x_2y_2/2-x_1y_1/2-x_2y_1         \\
         \hbox{}\quad \hbox{}
            =x_1y_2-x_2y_1        
    \end{array}$
    \hfil\hbox{}}
\end{center}
Kesamaan antara luas tersebut dan nilai determinan
\begin{equation*}
  \begin{vmat}
    x_1  &x_2  \\
    y_1  &y_2
  \end{vmat}
  =x_1y_2-x_2y_1
\end{equation*}
bukanlah kebetulan.
Definisi determinan memuat empat sifat yang kita ketahui
menghasilkan fungsi unik untuk setiap dimensi~$n$.
Kita akan menunjukkan bahwa sifat-sifat ini
merupakan postulat yang baik bagi fungsi
yang mengukur ukuran kotak dalam ruang-$n$.\index{ukuran}

Sebagai contoh, fungsi semacam itu
seharusnya memiliki sifat bahwa mengalikan salah satu vektor pembentuk kotak
dengan suatu skalar akan mengalikan ukurannya dengan skalar tersebut.
\begin{center}
  \includegraphics{det/mp/ch4.32}
  \qquad
  \includegraphics{det/mp/ch4.33}
\end{center}
Di sini ditampilkan $k=1.4$.
Di sebelah kanan, daerah yang diskalakan ulang digambar dengan garis utuh,
sedangkan daerah asal diarsir sebagai pembanding.
% The region formed by $k\vec{v}$ and~$\vec{w}$
% is bigger by a factor of \( k \)
% than the shaded region enclosed by $\vec{v}$ and~$\vec{w}$.
% That is,
% \( \size (k\vec{v},\vec{w})=k\cdot\size (\vec{v},\vec{w}) \) and

Artinya, wajar jika kita mengharapkan bahwa
$\size (\dots,k\vec{v},\dots)=k\cdot\size (\dots,\vec{v},\dots)$.
Tentu saja, syarat ini adalah salah satu syarat dalam
definisi determinan.

Sifat determinan lain yang seharusnya berlaku bagi setiap
fungsi yang mengukur ukuran kotak ialah bahwa
fungsi itu tidak terpengaruh oleh kombinasi baris.
Berikut adalah kotak sebelum dan sesudah pengombinasian
(skalar yang ditampilkan adalah $k=-0.35$).
\begin{center}
  \includegraphics{det/mp/ch4.34}
  \qquad
  \includegraphics{det/mp/ch4.35}
\end{center}   
Kotak yang dibentuk oleh $\vec{v}$ dan
$k\vec{v}+\vec{w}$ memiliki kemiringan yang berbeda
dari kotak semula, tetapi keduanya memiliki
alas dan tinggi yang sama, sehingga luasnya juga sama.
% (As before, the figure on the right has the 
% original region in shade for comparison.)
Jadi, kita mengharapkan bahwa ukuran tidak terpengaruh oleh operasi geser
$\size (\dots,\vec{v},\dots,\vec{w},\dots)
=\size (\dots,\vec{v},\dots,k\vec{v}+\vec{w},\dots)$.
Sekali lagi, ini adalah syarat determinan.

Kita mengharapkan bahwa kotak yang dibentuk oleh vektor-vektor satuan berukuran satu
\begin{center}
  \includegraphics{det/mp/ch4.36}
\end{center}
dan secara alami memperluasnya ke setiap ruang-$n$ sebagai
$\size(\vec{e}_1,\dots,\vec{e}_n)=1$.
% Again, that is a determinant property.

Syarat~(2) dalam definisi determinan
bersifat redundan, sebagaimana dinyatakan setelah definisi itu.
Dari bagian sebelumnya kita mengetahui bahwa
untuk setiap~$n$, determinan ada dan unik. Karena itu kita mengetahui
bahwa postulat-postulat untuk fungsi ukuran ini konsisten dan
bahwa kita tidak memerlukan postulat tambahan.
Dengan demikian, kita beralasan untuk
menafsirkan \( \det(\vec{v}_1,\dots,\vec{v}_n) \) sebagai ukuran
kotak yang dibentuk oleh vektor-vektor tersebut.
% (\textit{Comment.}
%   An even more basic approach, which also leads to the definition
%   below, is in \cite{Weston59}.)

\begin{remark}  \label{re:PropertyTwoGivesSign}
Meskipun syarat~(2) redundan, syarat itu menyoroti satu hal penting.
Perhatikan dua gambar berikut.
\begin{center} \small
  \begin{tabular}{c@{\hspace*{8em}}c}
    \includegraphics{det/mp/ch4.37}  
      &\includegraphics{det/mp/ch4.38}  \\[.25ex]
    \ $\begin{vmat}[r]
        4  &1   \\
        2  &3
      \end{vmat}=10$
      &\ $\begin{vmat}[r]
          1  &4   \\
          3  &2
        \end{vmat}=-10$
  \end{tabular}
\end{center}
Menukar kolom-kolom mengubah tandanya.
Di sebelah kiri, dengan mulai dari $\vec{u}$ dan mengikuti busur
di dalam sudut menuju $\vec{v}$ (yakni bergerak berlawanan arah jarum jam),
kita memperoleh ukuran positif.
Di sebelah kanan, dengan mulai dari $\vec{v}$ dan bergerak menuju~$\vec{u}$,
sehingga mengikuti busur searah jarum jam, kita memperoleh ukuran negatif.
Tanda yang dihasilkan oleh fungsi ukuran mencerminkan
\definend{orientasi}\index{kotak!orientasi}\index{orientasi}
atau \definend{arah}\index{kotak!arah}\index{arah} kotak tersebut.
(Hal yang sama terlihat jika kita membayangkan pengaruh perkalian skalar
dengan skalar negatif.)
% Although it is both interesting and important, we don't need the idea of
% orientation for the development below and so we will pass it by.
% (See \nearbyexercise{exer:BasisOrient}.)
\end{remark}

\begin{definition} \label{df:Volume}
%<*df:Volume>
\definend{Volume}\index{volume}\index{kotak!volume}
sebuah kotak adalah nilai mutlak determinan dari
matriks yang kolom-kolomnya berupa vektor-vektor tersebut.
%</df:Volume>
\end{definition}

\begin{example} \label{ex:VolParPiped}
Menurut rumus luas alas dikalikan tinggi,
volume paralelepiped ini adalah $12$.
Hasil ini sesuai dengan determinannya.
\begin{center}
   \parbox{2in}{\hbox{}\hfil\includegraphics{det/asy/ppiped.pdf}\hfil\hbox{}}  
  %  \parbox{2in}{\hbox{}\hfil\includegraphics{det/mp/ch4.39}\hfil\hbox{}}  
  \hspace{4.5em}
  $\begin{vmat}[r]
     2 &0 &-1\\
     0 &3 &0 \\
     2 &1 &1
  \end{vmat}=12$
\end{center}
Mengambil vektor-vektor dalam urutan berbeda mengubah tandanya, tetapi
tidak mengubah besar nilainya.
\begin{equation*}
  \begin{vmat}[r]
     0  &2 &-1 \\
     3  &0 &0 \\
     1  &2 &1
  \end{vmat}=-12  
\end{equation*}
\end{example}

% The next result describes some of the geometry of the linear
% functions that act on
% \( \Re^n \). 

\begin{theorem}\label{th:MatChVolByDetMat}
\index{ukuran}\index{transformasi!perubahan ukuran}
%<*th:MatChVolByDetMat>
Transformasi \( \map{t}{\Re^n}{\Re^n} \) mengubah ukuran semua kotak
dengan faktor yang sama; tepatnya, ukuran bayangan sebuah kotak
$\deter{t(S)}$ adalah $\deter{T}$ kali ukuran kotak $\deter{S}$,
di mana $T$ adalah matriks
yang merepresentasikan $t$ terhadap basis standar.

Dengan kata lain, determinan suatu hasil kali adalah
hasil kali determinan-determinannya, $\deter{TS}=\deter{T}\cdot\deter{S}$.
%</th:MatChVolByDetMat>
\end{theorem}

Kedua kalimat itu menyatakan hal yang sama, pertama dalam bahasa pemetaan lalu
dalam bahasa matriks.
Ini karena
$\deter{t(S)}=\deter{TS}$, sebab keduanya memberikan ukuran kotak yang merupakan
bayangan kotak satuan $\stdbasis_n$ di bawah komposisi $\composed{t}{s}$,
ketika pemetaan-pemetaan tersebut direpresentasikan terhadap basis standar.
Kita akan membuktikan kalimat kedua.

\begin{proof}
%<*pf:MatChVolByDetMat0>
Pertama, tinjau kasus ketika $T$ singular sehingga
tidak memiliki invers.
Perhatikan bahwa jika \( TS \) dapat diinverskan, maka terdapat
$M$ sedemikian sehingga \( (TS)M=I \), jadi
\( T(SM)=I \), sehingga \( T \) dapat diinverskan.
Kontraposisi pengamatan tersebut menyatakan bahwa
jika \( T \) tidak dapat diinverskan, maka \( TS \) juga tidak dapat diinverskan \Dash
jika $\deter{T}=0$, maka $\deter{TS}=0$.
%</pf:MatChVolByDetMat0>

%<*pf:MatChVolByDetMat1>
Sekarang tinjau kasus ketika $T$ nonsingular.
Setiap matriks nonsingular dapat difaktorkan menjadi hasil kali
matriks-matriks elementer $T=E_1E_2\cdots E_r$.
Untuk menyelesaikan argumen ini,
kita akan memverifikasi bahwa
\( \deter{ES}=\deter{E}\cdot\deter{S} \)
untuk semua matriks~$S$ dan matriks elementer~$E$.
Hasilnya kemudian mengikuti karena
$\deter{TS}=\deter{E_1\cdots E_rS}=\deter{E_1}\cdots\deter{E_r}\cdot\deter{S}
  =\deter{E_1\cdots E_r}\cdot\deter{S}=\deter{T}\cdot\deter{S}$.
%</pf:MatChVolByDetMat1>

%<*pf:MatChVolByDetMat2>
Ada tiga jenis matriks elementer.
Kita akan membahas kasus $M_i(k)$;
pemeriksaan untuk $P_{i,j}$ dan $C_{i,j}(k)$ serupa.
Matriks $M_i(k)S$ sama dengan $S$, kecuali bahwa baris~$i$ dikalikan dengan $k$.
Syarat ketiga fungsi determinan
kemudian memberikan $\deter{M_i(k)S}=k\cdot\deter{S}$.
Namun, $\deter{M_i(k)}=k$, sekali lagi menurut syarat ketiga, karena
$M_i(k)$ diperoleh dari identitas dengan mengalikan baris~$i$ dengan
$k$.
Jadi, \( \deter{ES}=\deter{E}\cdot\deter{S} \) berlaku untuk
$E=M_i(k)$.
%</pf:MatChVolByDetMat2>
\end{proof}


\begin{example}
Penerapan pemetaan $t$ yang direpresentasikan terhadap basis-basis standar oleh
\begin{equation*}
  \begin{mat}[r]
    1  &1  \\
   -2  &0
  \end{mat}
\end{equation*}
akan menggandakan ukuran kotak, misalnya dari ini
\begin{center}
    \parbox{1.5in}{\hbox{}\hfil\includegraphics{det/mp/ch4.40}\hfil\hbox{}}  
    \quad
    $\begin{vmat}[r]
      2  &1  \\
      1  &2
    \end{vmat}=3$
\end{center}
menjadi berikut ini
\begin{center}
    \parbox{1.5in}{\hbox{}\hfil\includegraphics{det/mp/ch4.41}\hfil\hbox{}}  
    \quad
    $\begin{vmat}[r]
        3  &3  \\
       -4  &-2
     \end{vmat}=6$
\end{center}
\end{example}


% Recall that determinants are not additive homomorphisms, that
% $\det(A+B)$ need not equal $\det(A)+\det(B)$.
% In contrast, the above theorem says that determinants are
% multiplicative homomorphisms:
% $\det(AB)$ equals $\det(A)\cdot \det(B)$.

\begin{corollary} \label{co:DeterminantOfInverseIsInverseOfDeterminant}
%<*co:DeterminantOfInverseIsInverseOfDeterminant>
Jika suatu matriks dapat diinverskan, maka determinan inversnya adalah
invers dari determinannya, $\deter{T^{-1}}=1/\deter{T}$.
%</co:DeterminantOfInverseIsInverseOfDeterminant>
\end{corollary}

\begin{proof}
%<*pf:DeterminantOfInverseIsInverseOfDeterminant>
$1=\deter{I}=\deter{TT^{-1}}=\deter{T}\cdot\deter{T^{-1}}$  
%</pf:DeterminantOfInverseIsInverseOfDeterminant>
\end{proof}


\begin{exercises}
  \item 
    Apakah
    \begin{equation*}
      \colvec[r]{4 \\ 1 \\ 2}
    \end{equation*}
    berada di dalam kotak yang dibentuk oleh ketiga vektor berikut?
    \begin{equation*}
      \colvec[r]{3 \\ 3 \\ 1}
      \quad
      \colvec[r]{2 \\ 6 \\ 1}
      \quad
      \colvec[r]{1 \\ 0 \\ 5}
    \end{equation*}
    \begin{answer}
      Menyelesaikan
      \begin{equation*}
        c_1\colvec[r]{3 \\ 3 \\ 1}
        +c_2\colvec[r]{2 \\ 6 \\ 1}
        +c_3\colvec[r]{1 \\ 0 \\ 5}
        =\colvec[r]{4 \\ 1 \\ 2}
      \end{equation*}
      memberikan solusi tunggal
      \( c_3=11/57 \), \( c_2=-40/57 \), dan \( c_1=99/57 \).
      Karena \( c_1>1 \), vektor tersebut tidak berada di dalam kotak.
    \end{answer}
  \recommended \item 
    Carilah volume daerah yang ditentukan oleh vektor-vektor tersebut.
    \begin{exparts}
      \partsitem $\sequence{\colvec[r]{1 \\ 3},\colvec[r]{-1 \\ 4}}$
      \partsitem $\sequence{\colvec[r]{2 \\ 1 \\ 0},\colvec[r]{3 \\ -2 \\ 4},
                              \colvec[r]{8 \\ -3 \\ 8}}$
      \partsitem $\sequence{\colvec[r]{1 \\ 2 \\ 0 \\ 1},
                             \colvec[r]{2 \\ 2 \\ 2 \\ 2},
                             \colvec[r]{-1 \\ 3 \\ 0 \\ 5},
                             \colvec[r]{0 \\ 1 \\ 0 \\ 7}}$
    \end{exparts}
    \begin{answer}
      Untuk setiap bagian, carilah determinannya lalu ambil nilai mutlaknya.
      \begin{exparts*}
        \partsitem $7$
        \partsitem $0$
         % sage: A = matrix(QQ, [[2, 3, 8], [1, -2, -3], [0, 4, 8]])
         % sage: A
         % [ 2  3  8]
         % [ 1 -2 -3]
         % [ 0  4  8]
         % sage: det(A)
         % 0
        \partsitem $58$
         % sage: A = matrix(QQ, [[1,2,-1,0], [2,2,3,1], [0,2,0,0], [1,2,5,7]])
         % sage: A
         % [ 1  2 -1  0]
         % [ 2  2  3  1]
         % [ 0  2  0  0]
         % [ 1  2  5  7]
         % sage: det(A)
         % -58
      \end{exparts*}
    \end{answer}
  \recommended \item 
    Pada gambar ini,
    persegi panjang di sebelah kiri ditentukan oleh titik-titik $(3,1)$ dan
    $(1,2)$.
    Terapkan matriks tersebut untuk memperoleh persegi panjang di sebelah kanan,
    yang ditentukan oleh $(7,1)$ dan $(4,2)$.
    Mengapa hal ini tidak bertentangan dengan
    \nearbytheorem{th:MatChVolByDetMat}?
    \begin{center}
      \begin{tabular}{c@{\hspace*{2em}}c@{\hspace*{2em}}c}
        \includegraphics{det/mp/ch4.43}
        &\raisebox{12pt}{\( \grstep{\bigl(\begin{smallmatrix}
                                        2  &1 \\
                                        0  &1
                                     \end{smallmatrix}\bigr)} \)}
        &\includegraphics{det/mp/ch4.44}                                 \\
        luasnya $2$
        &determinannya $2$
        &luasnya $3$
      \end{tabular}
    \end{center}
    \begin{answer}
      Gambar itu sengaja dibuat menyesatkan.
      Gambar di sebelah kiri bukan kotak yang dibentuk oleh dua vektor.
      Jika kita menggesernya ke titik asal, gambar itu menjadi kotak yang
      dibentuk oleh barisan berikut.
      \begin{equation*}
        \sequence{
          \colvec[r]{0 \\ 1},
          \colvec[r]{2 \\ 0}
       }
      \end{equation*}
      Kemudian, bayangan di bawah aksi matriks tersebut adalah kotak yang
      dibentuk oleh barisan berikut.
      \begin{equation*}
        \sequence{
          \colvec[r]{1 \\ 1},
          \colvec[r]{4 \\ 0}
         }
      \end{equation*}
      yang luasnya $4$.
     \end{answer}
  \recommended \item 
    Carilah volume daerah ini.
    \begin{center}
      \includegraphics{det/mp/ch4.42}  
    \end{center}
    \begin{answer}
      Geser paralelepiped itu agar bermula di titik asal,
      sehingga menjadi kotak yang dibentuk oleh
      \begin{equation*}
        \sequence{
          \colvec[r]{3 \\ 0},
          \colvec[r]{2 \\ 1}
        }      
      \end{equation*}
      dan sekarang nilai mutlak determinan ini
      mudah dihitung sebagai $3$.
      \begin{equation*}
        \begin{vmat}[r]
          3  &2  \\
          0  &1
        \end{vmat}=3
      \end{equation*}
     \end{answer}
  \recommended \item 
    Misalkan \( \deter{A}=3 \).
    Dengan faktor berapa masing-masing matriks berikut mengubah volume?
    \begin{exparts*}
      \partsitem \( A \)
      \partsitem \( A^2 \)
      \partsitem \( A^{-2} \)
    \end{exparts*}
    \begin{answer}
     \begin{exparts*}
        \partsitem \( 3 \)
        \partsitem \( 9 \)
        \partsitem $1/9$
      \end{exparts*}  
    \end{answer}
  \recommended \item
    Tinjau transformasi linear~$t$ pada~$\Re^3$
    yang direpresentasikan terhadap
    basis-basis standar oleh matriks ini.
    \begin{equation*}
      \begin{mat}
        1 &0 &-1 \\
        3 &1 &1 \\
       -1 &0 &3
      \end{mat}
    \end{equation*}
    \begin{exparts}
      \partsitem Hitung determinan matriks tersebut.
        Apakah transformasi itu mempertahankan atau membalik orientasi?
      \partsitem Carilah ukuran kotak yang ditentukan oleh vektor-vektor berikut.
        Apa orientasinya?
        \begin{equation*}
          \colvec{1 \\ -1 \\ 2}
          \quad
          \colvec{2 \\ 0 \\ -1}
          \quad
          \colvec{1 \\ 1 \\ 0}
        \end{equation*}
      \partsitem Carilah bayangan di bawah $t$ dari vektor-vektor pada bagian
        sebelumnya, lalu carilah ukuran kotak yang dibentuknya.
        Apa orientasinya?
    \end{exparts}
    \begin{answer}
      \begin{exparts}
        \partsitem
          Metode Gauss
          \begin{equation*}
          \grstep[\rho_1+\rho_3]{-3\rho_1+\rho_2}
          \begin{mat}
            1 &0   &-1    \\
            0 &1   &4     \\
            0 &0   &2
          \end{mat}
          \end{equation*}
          memberikan determinan~$+2$.
          Tandanya positif, sehingga transformasi itu mempertahankan orientasi.
        \partsitem
          Ukuran kotak adalah nilai determinan berikut.
          \begin{equation*}
            \begin{vmat}
              1 &2  &1 \\
             -1 &0  &1 \\
              2 &-1 &0
            \end{vmat}
            =+6
          \end{equation*}
          Orientasinya positif.
        \partsitem
          Karena transformasi ini direpresentasikan oleh matriks yang diberikan
          terhadap basis-basis standar, dan terhadap basis standar vektor-vektor
          merepresentasikan dirinya sendiri, untuk mencari bayangan vektor-vektor
          di bawah transformasi tersebut cukup kalikan vektor-vektor itu dari
          kiri dengan matriks.
          \begin{equation*}
            \colvec{1 \\ -1 \\ 2}\mapsto\colvec{-1 \\ 4 \\ 5}
            \qquad
            \colvec{2 \\ 0 \\ -1}\mapsto\colvec{3 \\ 5 \\ -5}
            \qquad
            \colvec{1 \\ 1 \\ 0}\mapsto\colvec{1 \\ 4 \\ -1}
          \end{equation*}
          Kemudian hitung ukuran kotak yang dihasilkan.
          \begin{equation*}
            \begin{vmat}
              -1 &3  &1 \\
               4 &5  &4 \\
               5 &-5 &-1
            \end{vmat}
            =+12
          \end{equation*}
          Kotak awal berorientasi positif, transformasinya mempertahankan
          orientasi (karena determinan matriksnya positif), dan kotak akhir
          juga berorientasi positif.
      \end{exparts}
    \end{answer}
  \item 
    Dengan faktor berapa setiap transformasi mengubah ukuran
    kotak?
    \begin{exparts*}
      \partsitem $\colvec{x \\ y}\mapsto\colvec{2x \\ 3y}$
      \partsitem $\colvec{x \\ y}\mapsto\colvec{3x-y \\ -2x+y}$
      \partsitem $\colvec{x \\ y \\ z}\mapsto\colvec{x-y \\ x+y+z \\ y-2z}$
    \end{exparts*}
    \begin{answer}
      Nyatakan setiap transformasi terhadap basis-basis standar
      lalu carilah determinannya.
      \begin{exparts*}
        \partsitem $6$
        \partsitem $-1$
        \partsitem $-5$
      \end{exparts*}
    \end{answer}
  \item 
    Berapakah luas bayangan persegi panjang
    \( [2..4]\times [2..5] \) di bawah aksi
    matriks berikut?
    \begin{equation*}
       \begin{mat}[r]
         2  &3  \\
         4  &-1
       \end{mat}
    \end{equation*}
    \begin{answer}
      Luas awal adalah \( 6 \), dan matriks tersebut mengubah ukuran dengan
      faktor \( -14 \).
      Jadi, luas bayangannya adalah \( 84 \).
    \end{answer}
  \item
     Jika \( \map{t}{\Re^3}{\Re^3} \) mengubah volume dengan faktor \( 7 \)
     dan \( \map{s}{\Re^3}{\Re^3} \) mengubah volume dengan faktor \( 3/2 \),
     maka dengan faktor berapa komposisi keduanya mengubah volume?
     \begin{answer}
        Dengan faktor \( 21/2 \).
     \end{answer}
  \item 
    Dalam hal apa definisi kotak berbeda dari
    definisi rentang?
    \begin{answer}
      Untuk kotak, kita mengambil barisan vektor (sebagaimana dijelaskan
      dalam catatan, urutan vektor penting), sedangkan untuk rentang kita
      mengambil himpunan vektor.
      Selain itu, untuk kotak yang merupakan himpunan bagian $\Re^n$ harus ada
      $n$ vektor; tentu saja, suatu rentang dapat memiliki berapa pun banyaknya
      vektor.
      Terakhir, untuk kotak, koefisien $t_1$,~\ldots, $t_n$ berada dalam
      interval $[0..1]$, sedangkan untuk rentang, koefisiennya bebas mengambil
      nilai di seluruh $\Re$.
    \end{answer}
  \item 
    Apakah \( \deter{TS}=\deter{ST} \)?
    \( \deter{T(SP)}=\deter{(TS)P} \)?
    \begin{answer}
      Ya, untuk keduanya.
      Sebagai contoh, persamaan pertama adalah \( \deter{TS}=\deter{T}\cdot\deter{S}=
                     \deter{S}\cdot\deter{T}=\deter{ST} \).  
    \end{answer}
  \item Tunjukkan bahwa tidak ada matriks $\nbyn{2}$ $A$ dan~$B$
   yang memenuhi persamaan berikut. % http://math.stackexchange.com/questions/827262/need-help-with-a-linear-algebra-proof
   \begin{equation*}
     AB=\begin{mat}[r]
       1 &-1 \\
       2  &0
     \end{mat}
     \quad
     BA=\begin{mat}[r]
       2 &1 \\
       1  &1
     \end{mat}
   \end{equation*}
   \begin{answer}  % due to math.stackexchange.com user dgrasines517
     Karena $\deter{AB}=\deter{A}\cdot\deter{B}=\deter{BA}$, sedangkan kedua
     matriks tersebut memiliki determinan yang berbeda.
   \end{answer}
  \item 
   \begin{exparts}
     \partsitem Misalkan \( \deter{A}=3 \) dan \( \deter{B}=2 \).
        Carilah \( \deter{A^2\cdot \trans{B}\cdot B^{-2}\cdot \trans{A} } \).
     \partsitem Andaikan \( \deter{A}=0 \).
        Buktikan bahwa \( \deter{6A^3+5A^2+2A}=0 \).
    \end{exparts}
    \begin{answer}
      \begin{exparts}
        \partsitem Jika terdefinisi, nilainya adalah
           \( (3^2)\cdot (2)\cdot (2^{-2})\cdot (3) \).
        \partsitem \( \deter{6A^3+5A^2+2A}=\deter{A}\cdot\deter{6A^2+5A+2I} \).
      \end{exparts}  
    \end{answer}
  \recommended \item
    Misalkan \( T \) adalah matriks yang merepresentasikan (terhadap basis-basis
    standar) pemetaan yang memutar vektor-vektor bidang sejauh
    \( \theta \) radian berlawanan arah jarum jam.
    Dengan faktor berapa \( T \) mengubah ukuran?
    \begin{answer}
       \(\begin{vmat}
                \cos\theta  &-\sin\theta  \\
                \sin\theta  &\cos\theta
              \end{vmat}=1 \)  
    \end{answer}
  \recommended \item
    Apakah transformasi \( \map{t}{\Re^2}{\Re^2} \) yang mempertahankan luas
    harus pula mempertahankan panjang?
    \begin{answer}
      Tidak. Sebagai contoh, determinan
      \begin{equation*}
        T=\begin{mat}[r]
          2  &0  \\
          0  &1/2
        \end{mat}
      \end{equation*}
      adalah \( 1 \), sehingga matriks itu mempertahankan luas, tetapi vektor
      \( T\vec{e}_1 \) memiliki panjang \( 2 \).
    \end{answer}
  \item
    Berapakah volume paralelepiped di \( \Re^3 \) yang dibatasi oleh
    himpunan bergantung linear?
    \begin{answer}
       Volumenya nol.
    \end{answer}
  \recommended \item
    Carilah luas segitiga di \( \Re^3 \) yang titik-titik sudutnya
    \( (1,2,1) \), \( (3,-1,4) \), dan \( (2,2,2) \).
    (Yang diminta adalah luas, bukan volume.
    Segitiga itu menentukan sebuah bidang; berapakah luas segitiga pada
    bidang tersebut?)
    \begin{answer}
      Dua dari tiga sisi segitiga dibentuk oleh vektor-vektor berikut.
      \begin{equation*}
        \colvec[r]{2 \\ 2 \\ 2}-\colvec[r]{1 \\ 2 \\ 1}=\colvec[r]{1 \\ 0 \\ 1}
        \qquad
        \colvec[r]{3 \\ -1 \\ 4}-\colvec[r]{1 \\ 2 \\ 1}=\colvec[r]{2 \\ -3 \\ 3}
      \end{equation*}
      Salah satu cara mencari luas segitiga ini adalah mencari setengah volume
      paralelepiped yang dibentuk oleh kedua vektor tersebut dan sebuah vektor
      panjang satu yang ortogonal terhadap keduanya.
      Untuk mencari keluarga vektor yang ortogonal terhadap keduanya, kita dapat
      mulai dengan dua relasi
      \begin{equation*}
        \colvec[r]{1 \\ 0 \\ 1}
        \dotprod\colvec{x \\ y \\ z}
        =0
        \qquad
        \colvec[r]{2 \\ -3 \\ 3}
        \dotprod\colvec{x \\ y \\ z}
        =0
      \end{equation*}
      lalu menyelesaikan sistem
      \begin{equation*}
        \begin{linsys}{3}
          x  &   &   &+  &z  &=  &0  \\
         2x  &-  &3y &+  &3z &=  &0  
        \end{linsys}
        \grstep{-2\rho_1+\rho_2}
        \begin{linsys}{3}
          x  &   &   &+  &z  &=  &0  \\
             &   &-3y&+  &z  &=  &0  
        \end{linsys}
      \end{equation*}
      untuk memperoleh himpunan solusi berikut.
      \begin{equation*}
        \set{\colvec[r]{-1 \\ 1/3 \\ 1}z\suchthat z\in\Re}
      \end{equation*}
      Berikut adalah solusi dengan panjang satu.
      \begin{equation*}
        \frac{1}{\sqrt{19/9}}\cdot\colvec[r]{-1 \\ 1/3 \\ 1}
        =\colvec[r]{-3/\sqrt{19} \\ 1/\sqrt{19} \\ 3/\sqrt{19}}
      \end{equation*}
      Jadi, luas segitiga adalah setengah dari nilai mutlak
      determinan berikut.
      \begin{equation*}
        \begin{vmat}[r]
             1  &2   &-3/\sqrt{19}   \\
             0  &-3  &1/\sqrt{19}   \\
             1  &3   &3/\sqrt{19}
        \end{vmat}
        =-19/\sqrt{19}
      \end{equation*} 
      Setengah dari nilai mutlak tersebut adalah $\sqrt{19}/2$.
    \end{answer}
  \item \label{exer:DetProdEqProdDetsFcn}
    Pembuktian alternatif untuk \nearbytheorem{th:MatChVolByDetMat}
    menggunakan definisi fungsi determinan.
    \begin{exparts}
      \partsitem Perhatikan bahwa vektor-vektor pembentuk
        $S$ membentuk himpunan bergantung linear jika dan hanya jika
        $\deter{S}=0$, lalu periksa bahwa hasil tersebut berlaku dalam kasus ini.
      \partsitem Untuk kasus $\deter{S}\neq 0$, guna menunjukkan bahwa
        $\deter{TS}/\deter{S}=\deter{T}$ bagi semua transformasi, tinjau fungsi
        \( \map{d}{\matspace_{\nbyn{n}}}{\Re} \) yang diberikan oleh
        \( T\mapsto \deter{TS}/\deter{S} \).
        Tunjukkan bahwa $d$ memiliki sifat pertama determinan.
      \partsitem Tunjukkan bahwa $d$ memiliki tiga sifat lainnya dari
        fungsi determinan.
      \partsitem Simpulkan bahwa $\deter{TS}=\deter{T}\cdot\deter{S}$.
    \end{exparts}
    \begin{answer}
      \begin{exparts}
        \partsitem Karena bayangan himpunan bergantung linear juga
          bergantung linear, jika vektor-vektor pembentuk $S$ membentuk
          himpunan bergantung linear sehingga $\deter{S}=0$, maka
          vektor-vektor pembentuk $t(S)$ juga membentuk himpunan bergantung
          linear sehingga $\deter{TS}=0$; dalam kasus ini persamaan tersebut berlaku.
        \partsitem Kita harus memeriksa bahwa jika
          $T\smash[b]{\grstep{k\rho_i+\rho_j}}\hat{T}$, maka
          $d(T)=\deter{TS}/\deter{S}=\deter{\hat{T}S}/\deter{S}=d(\hat{T})$.
          Kita dapat melakukannya dengan memeriksa bahwa mengombinasikan baris
          terlebih dahulu lalu mengalikan untuk memperoleh \( \hat{T}S \)
          memberikan hasil yang sama dengan mengalikan terlebih dahulu untuk
          memperoleh \( TS \) lalu mengombinasikan baris (karena determinan
          \( \deter{TS} \) tidak terpengaruh oleh pengombinasian baris, sehingga
          kita memperoleh \( \deter{\hat{T}S}=\deter{TS} \), dan karenanya
          \( d(\hat{T})=d(T) \)).
          Pemeriksaannya sebagai berikut:~setelah menambahkan
          \( k \) kali baris~\( i \) dari \( TS \) ke
          baris~$j$ dari \( TS \), entri \( j,p \) adalah
          \( (kt_{i,1}+t_{j,1})s_{1,p}+\dots+(kt_{i,r}+t_{j,r})s_{r,p} \),
          yang merupakan entri \( j,p \) dari \( \hat{T}S \).
        \partsitem Untuk sifat kedua, kita hanya perlu memeriksa bahwa melakukan
          pertukaran
          $T\smash[b]{\grstep{\rho_i\swap\rho_j}}\hat{T}$
          lalu mengalikan untuk memperoleh \( \hat{T}S \) memberikan hasil yang
          sama dengan mengalikan \( T \) dengan \( S \) terlebih dahulu lalu
          menukar baris (karena determinan \( \deter{TS} \) berubah tanda ketika
          baris ditukar, sehingga kita memperoleh
          \( \deter{\hat{T}S}=-\deter{TS} \), dan jadi
          \( d(\hat{T})=-d(T) \)).
          Pemeriksaan ini berlangsung sama seperti untuk sifat pertama.

          Untuk sifat ketiga, kita hanya perlu menunjukkan bahwa melakukan
          $T\smash[b]{\grstep{k\rho_i}}\hat{T}$ 
          lalu menghitung \( \hat{T}S \) memberikan hasil yang sama dengan
          menghitung \( TS \) terlebih dahulu lalu melakukan perkalian skalar
          (karena determinan \( \deter{TS} \) diskalakan ulang dengan \( k \),
          kita akan memperoleh \( \deter{\hat{T}S}=k\deter{TS} \), sehingga
          \( d(\hat{T})=k\,d(T) \)).
          Di sini pun, argumennya berlangsung seperti di atas.

          Sifat keempat, yakni jika $T$ adalah $I$ maka hasilnya $1$,
          sudah jelas.
        \partsitem Fungsi determinan bersifat unik, sehingga
          \( \deter{TS}/\deter{S}=d(T)=\deter{T} \),
          dan jadi $\deter{TS}=\deter{T}\deter{S}$.
      \end{exparts}
    \end{answer}
%  \item  
%    Use the fact that
%    \( \deter{TS}=\deter{T}\,\deter{S} \)
%    to prove that if \( \phi \) and \( \sigma \) are \( n \)-permutations
%    then \( \sgn(\phi\sigma)=\sgn(\phi)\sgn(\sigma) \).
%    \cite{HoffmanKunze}
%    \begin{answer}
%       Take \( T=P_\phi \) and \( S=P_\sigma \).
%      Note that \( TS=P_{\phi\sigma} \) and so
%      \( \deter{P_{\phi\sigma}}=\deter{P_\phi}\cdot\deter{P_\sigma} \).  
%    \end{answer}
  \item 
    Berikan sebuah matriks bukan identitas dengan sifat
    \( \trans{A}=A^{-1} \).
    Tunjukkan bahwa jika \( \trans{A}=A^{-1} \), maka \( \deter{A}=\pm 1 \).
    Apakah konversnya berlaku?
    \begin{answer}
      Setiap matriks permutasi memiliki sifat bahwa transpos matriks tersebut
      adalah inversnya.

      Untuk implikasi tersebut, kita mengetahui bahwa
      \( \deter{\trans{A}}=\deter{A} \).
      Maka \( 1=\deter{A\cdot A^{-1}}=\deter{A\cdot\trans{A}}
               =\deter{A}\cdot\deter{\trans{A}}=\deter{A}^2 \).

      Konversnya tidak berlaku; berikut adalah sebuah contoh.
      \begin{equation*}
        \begin{mat}[r]
          3  &1  \\
          2  &1
        \end{mat}
      \end{equation*}
    \end{answer}
  \item 
    Sifat aljabar determinan bahwa mengeluarkan suatu skalar sebagai faktor
    dari satu baris akan mengalikan determinan dengan skalar tersebut
    menunjukkan bahwa jika \( H \) berukuran \( \nbyn{3} \), determinan
    \( cH \) adalah \( c^3 \) kali determinan \( H \).
    Jelaskan hal ini secara geometris, yakni dengan menggunakan
    \nearbytheorem{th:MatChVolByDetMat}.
    (Pengamatan bahwa memperbesar ukuran linear suatu benda tiga dimensi dengan
    faktor $c$ akan memperbesar volumenya dengan faktor $c^3$, sementara luas
    permukaannya hanya bertambah dengan faktor yang sebanding dengan $c^2$,
    disebut \definend{hukum kuadrat--kubik}~\cite{Wikipedia}.)
    \begin{answer}
      Ketika sisi-sisi kotak \( c \) kali lebih panjang, volumenya memuat
      \( c^3 \) kali lebih banyak satuan kubik.
    \end{answer}
  \item 
    Kita mengatakan bahwa matriks $H$ dan $G$
    \definend{serupa}\index{serupa}\index{matriks!serupa}
    jika terdapat matriks nonsingular $P$ sedemikian sehingga $H=P^{-1}GP$
    (kita akan mempelajari relasi ini dalam Bab Lima).
    Tunjukkan bahwa matriks-matriks serupa memiliki determinan yang sama.
    \begin{answer}
       Jika \( H=P^{-1}GP \),
       maka \( \deter{H}=\deter{P^{-1}}\deter{G}\deter{P}
        =\deter{P^{-1}}\deter{P}\deter{G}=\deter{P^{-1}P}\deter{G}
        =\deter{G} \).  
    \end{answer}
  \item  \label{exer:BasisOrient}
    Biasanya kita merepresentasikan vektor-vektor di \( \Re^2 \) terhadap
    basis standar sehingga vektor-vektor di kuadran pertama memiliki kedua
    koordinat positif.
    \begin{center}
      \parbox{.75in}{\hbox{}\hfil\includegraphics{det/mp/ch4.45}\hfil\hbox{}}
      \qquad
      \( \rep{\vec{v}}{\stdbasis_2}=\colvec[r]{+3 \\ +2} \)
    \end{center}
    Dengan bergerak berlawanan arah jarum jam mengelilingi titik asal, kita
    melewati empat daerah secara bersiklus:
    {\scriptsize
    \begin{equation*}
       \cdots
       \;\longrightarrow\colvec{+ \\ +}
       \;\longrightarrow\colvec{- \\ +}
       \;\longrightarrow\colvec{- \\ -}
       \;\longrightarrow\colvec{+ \\ -}
       \;\longrightarrow\cdots\,.
    \end{equation*} }
    Menggunakan basis berikut
    \begin{center}
      \( B=\sequence{\colvec[r]{0 \\ 1},\colvec[r]{-1 \\ 0}} \)
      \qquad
      \parbox{.75in}{\hbox{}\hfil\includegraphics{det/mp/ch4.46}\hfil\hbox{}}
    \end{center}
    menghasilkan siklus berlawanan arah jarum jam yang sama.
    Kita mengatakan bahwa kedua basis ini memiliki \emph{orientasi} yang sama.\index{orientasi}
    \begin{exparts}
      \partsitem Mengapa keduanya menghasilkan siklus yang sama?
      \partsitem Konfigurasi vektor satuan lain apa pada sumbu-sumbu yang
        menghasilkan siklus yang sama?
      \partsitem Carilah determinan matriks-matriks yang dibentuk dari
        basis-basis (terurut) tersebut.
      \partsitem Siklus lain apa yang mungkin terjadi ketika bergerak berlawanan
        arah jarum jam, dan berapakah determinan yang terkait?
      \partsitem Apa yang terjadi di \( \Re^1 \)?
      \partsitem Apa yang terjadi di \( \Re^3 \)?
    \end{exparts}
    Pembahasan orientasi yang menarik bagi pembaca umum
    terdapat dalam \cite{Gardner}.
    \begin{answer}
      \begin{exparts}
        \partsitem Basis baru adalah basis lama yang diputar sebesar \( \pi/2 \).
        \partsitem 
          $
             \sequence{\colvec[r]{-1 \\ 0},
                       \colvec[r]{0 \\ -1}}
          $, $
             \sequence{\colvec[r]{0 \\ -1},
                       \colvec[r]{1 \\ 0}}
          $
        \partsitem Dalam setiap kasus, determinannya adalah \( +1 \)
          (kita mengatakan bahwa basis-basis ini memiliki
          \definend{orientasi positif}).
        \partsitem Karena hanya satu tanda dapat berubah pada satu waktu,
          satu-satunya siklus lain yang mungkin adalah
          \begin{equation*}
             \cdots
             \;\longrightarrow\colvec{+ \\ +}
             \;\longrightarrow\colvec{+ \\ -}
             \;\longrightarrow\colvec{- \\ -}
             \;\longrightarrow\colvec{- \\ +}
             \;\longrightarrow\cdots\,.
          \end{equation*}
          Di sini setiap determinan yang terkait adalah \( -1 \)
          (kita mengatakan bahwa basis-basis tersebut memiliki
          \definend{orientasi negatif}).
        \partsitem Ada satu basis berorientasi positif \( \sequence{(1)} \)
          dan satu basis berorientasi negatif \( \sequence{(-1)} \).
        \partsitem Ada \( 48 \) basis (terdapat \( 6 \) pilihan setengah-sumbu
          untuk vektor satuan pertama, \( 4 \) untuk vektor kedua, dan
          \( 2 \) untuk vektor terakhir).
          Separuhnya berorientasi positif seperti basis standar di sebelah kiri
          di bawah, dan separuhnya berorientasi negatif seperti basis di sebelah
          kanan
         \begin{center}
           \includegraphics{det/mp/ch4.47}
           \hspace*{4em}
           \includegraphics{det/mp/ch4.48}
          \end{center}
          Di \( \Re^3 \), orientasi positif kadang-kadang disebut
          `orientasi tangan kanan' karena jika seseorang menempatkan tangan
          kanannya dengan jari-jari melengkung dari \( \vec{e}_1 \) menuju
          \( \vec{e}_2 \), maka ibu jarinya akan menunjuk searah dengan
          \( \vec{e}_3 \).
      \end{exparts}  
    \end{answer}
%  \item 
%    A region of \( \Re^n \) is \definend{convex}\index{convex region} 
%    if for any two points
%    connecting that region, the line segment joining them lies entirely
%    inside the region (the inside of a sphere is convex, while the skin of a
%    sphere or a horseshoe is not).
%    Prove that boxes are convex.
%    \begin{answer}
%      Let \( \vec{p}=p_1\vec{v}_1+\dots +p_n\vec{v}_n \) and
%      \( \vec{q}=q_1\vec{v}_1+\dots +q_n\vec{v}_n \) be two vectors from a box
%      so that \( p_1,\dots,\,p_n,q_1,\dots,\,q_n\in [0..1] \).
%      The line segment between them is this.
%      \begin{equation*}
%        \set{t\cdot\vec{p}+(1-t)\cdot\vec{q}
%              \suchthat t\in [0..1]}
%        =\set{(tp_1+(1-t)q_1)\cdot\vec{v}_1+\dots+(tp_n+(1-t)q_n)\cdot\vec{v}_n
%              \suchthat t\in [0..1] }
%      \end{equation*}
%      Showing that each member of that set is in the box is routine.  
%    \end{answer}
  \item \label{exer:DetProdEqProdDetsPerms}
    \textit{Soal ini menggunakan materi dari subbagian opsional
      Keberadaan Fungsi Determinan.}
    Buktikan \nearbytheorem{th:MatChVolByDetMat} dengan menggunakan
    rumus ekspansi permutasi untuk determinan.
    \begin{answer}
      Kita akan membandingkan \( \det(\vec{s}_1,\dots,\vec{s}_n) \) dengan
      \( \det(t(\vec{s}_1),\dots,t(\vec{s}_n)) \) untuk menunjukkan bahwa
      yang kedua berbeda dari yang pertama dengan faktor $\deter{T}$.
      Kita merepresentasikan vektor-vektor \( \vec{s}\, \) terhadap basis-basis standar
      \begin{equation*}
        \rep{\vec{s}_i}{\stdbasis_n}=
           \colvec{s_{1,i} \\ s_{2,i} \\ \vdots \\ s_{n,i}}
      \end{equation*}
      lalu kita merepresentasikan penerapan pemetaan tersebut dengan
      perkalian matriks-vektor
      \begin{align*}
        \rep{\,t(\vec{s}_i)\,}{\stdbasis_n}
         &=\generalmatrix{t}{n}{n}
           \colvec{s_{1,i} \\ s_{2,i} \\ \vdots \\ s_{n,i}}     \\
         &=s_{1,i}\colvec{t_{1,1} \\ t_{2,1} \\ \vdots \\ t_{n,1}}
          +s_{2,i}\colvec{t_{1,2} \\ t_{2,2} \\ \vdots \\ t_{n,2}}
          +\dots
          +s_{n,i}\colvec{t_{1,n} \\ t_{2,n} \\ \vdots \\ t_{n,n}} \\
         &=s_{1,i}\vec{t}_1+s_{2,i}\vec{t}_2+\dots+s_{n,i}\vec{t}_n
      \end{align*}
      di mana \( \vec{t}_i \) adalah kolom~$i$ dari \( T \).
      Maka $\det(t(\vec{s}_1),\,\dots,\,t(\vec{s}_n))$ sama dengan
      $
        \det(s_{1,1}\vec{t}_1\!+\!s_{2,1}\vec{t}_2\!
                +\!\dots\!+\!s_{n,1}\vec{t}_n,\,
             \dots,\,
             s_{1,n}\vec{t}_1\!+\!s_{2,n}\vec{t}_2
                \!+\!\dots\!+\!s_{n,n}\vec{t}_n)
     $.

     Seperti dalam penurunan rumus ekspansi permutasi, kita
     menerapkan multilinearitas, pertama-tama dengan memecah menurut
     penjumlahan pada argumen pertama
     \begin{multline*}
         \det(s_{1,1}\vec{t}_1,\,
            \dots,\,
            s_{1,n}\vec{t}_1+s_{2,n}\vec{t}_2+\dots+s_{n,n}\vec{t}_n)  \\ 
        +\cdots{}                                             
        +\det(s_{n,1}\vec{t}_n,\,
           \ldots,\,
            s_{1,n}\vec{t}_1+s_{2,n}\vec{t}_2+\dots+s_{n,n}\vec{t}_n)
     \end{multline*}
     lalu memecah masing-masing dari $n$ suku tersebut menurut penjumlahan
     pada argumen kedua, dan seterusnya.
     Seperti dalam penurunan ekspansi permutasi, pada akhirnya kita memperoleh
     \( n^n \) suku determinan, masing-masing berbentuk
     $\det(s_{i_1,1}\vec{t}_{i_1},s_{i_2,2}\vec{t}_{i_2},
            \,\dots,\,
            s_{i_n,n}\vec{t}_{i_n})$.
     Keluarkan setiap $s_{i,j}$ sebagai faktor
     $=
       s_{i_1,1}s_{i_2,2}\dots s_{i_n,n}
        \cdot\det(\vec{t}_{i_1},\vec{t}_{i_2},
       \,\dots,\,
       \vec{t}_{i_n})
      $.

      Seperti dalam penurunan ekspansi permutasi,
      apabila dua indeks di antara $i_1$, \ldots, $i_n$ sama, determinannya
      memiliki dua argumen yang sama dan bernilai $0$.
      Jadi, kita hanya perlu meninjau kasus ketika $i_1$, \ldots, $i_n$
      membentuk permutasi dari bilangan $1$, \ldots, $n$.
      Dengan demikian, kita memperoleh
      \begin{equation*}
        \det(t(\vec{s}_1),\dots,t(\vec{s}_n))=
          \sum_{\text{permutasi\ } \phi}
            s_{\phi(1),1}\dots s_{\phi(n),n}
            \det(\vec{t}_{\phi(1)},\dots,\vec{t}_{\phi(n)}).
      \end{equation*}
      Tukarkan kolom-kolom dalam
      $\det(\vec{t}_{\phi(1)},\ldots,\vec{t}_{\phi(n)})$ untuk memperoleh
      kembali matriks \( T \); ini mengubah tanda dengan faktor $\sgn{\phi}$.
      Kemudian keluarkan determinan $T$ sebagai faktor.
      \begin{equation*}
        =\sum_\phi
          s_{\phi(1),1}\dots s_{\phi(n),n}
            \det(\vec{t}_1,\dots,\vec{t}_n)\cdot\sgn{\phi}
        =\det(T)\sum_\phi
          s_{\phi(1),1}\dots s_{\phi(n),n}\cdot\sgn{\phi}.
      \end{equation*}
      Seperti dalam pembuktian bahwa determinan suatu matriks sama dengan
      determinan transposnya, kita mempertukarkan faktor-faktor $s$ agar
      tercantum menurut nomor baris yang menaik, bukan menurut nomor kolom
      yang menaik (dan kita mengganti $\sgn(\phi)$ dengan
      $\sgn(\phi^{-1})$).
      \begin{equation*}
        =\det(T)\sum_\phi
          s_{1,\phi^{-1}(1)}\dots s_{n,\phi^{-1}(n)}\cdot\sgn{\phi^{-1}}  
        =\det(T)\det(\vec{s}_1,\vec{s}_2,\dots,\vec{s}_n)
      \end{equation*}
    \end{answer}
%  \item 
%    Suppose that \( \map{f}{\matspace_{\nbyn{n}}}{\Re} \) is a non-constant
%    function with the property that \( f(GH)=f(G)f(H) \).
%    \begin{exparts}
%      \partsitem Show that \( f \) sends the identity to \( 1 \).
%      \partsitem Show that \( f \) maps the elementary matrix 
%        \( C_{i,j}(k) \) to \( 1 \)
%        (this matrix results from performing \( k\rho_i+\rho_j \) to the
%        identity).
%      \partsitem Show that 
%        \( f \) maps a row swap matrix to \( +1 \) or \( -1 \).
%    \end{exparts}
%    \begin{answer}
%      \begin{exparts}
%        \partsitem We have \( f(IH)=f(I)f(H) \) and \( f(IH)=f(H) \).
%        \partsitem
%        \partsitem A row swap matrix has the property that when 
%          done twice it equals
%          the identity.
%          But \( f(R)f(R)=f(RR)=f(I)=1 \) implies that \( f(R)=\pm 1 \).
%      \end{exparts} 
%     \end{answer}
  \recommended \item
    \begin{exparts}
      \partsitem Tunjukkan bahwa bentuk berikut memberikan
        persamaan garis di \( \Re^2 \) yang melalui
        \( (x_2,y_2) \) dan \( (x_3,y_3) \).
        \begin{equation*}
          \begin{vmat}
            x    &x_2 &x_3  \\
            y    &y_2 &y_3  \\
            1    &1   &1
          \end{vmat}=0
        \end{equation*}
      \partsitem \cite{Monthly55p249}
        Buktikan bahwa luas segitiga dengan titik-titik sudut \( (x_1,y_1) \),
        \( (x_2,y_2) \), dan \( (x_3,y_3) \) adalah
        \begin{equation*}
          \frac{1}{2}\left|
          \begin{vmat}
            x_1  &x_2 &x_3  \\
            y_1  &y_2 &y_3  \\
            1    &1   &1
          \end{vmat}\right|.
        \end{equation*}
      \partsitem \cite{MathMag73p286}
        Buktikan bahwa segitiga dengan titik-titik sudut \( (x_1,y_1) \),
        \( (x_2,y_2) \), dan \( (x_3,y_3) \), yang semua koordinatnya berupa
        bilangan bulat, memiliki luas \( N \) atau \( N/2 \) untuk suatu
        bilangan bulat positif \( N \).
    \end{exparts}
    \begin{answer}
      \begin{exparts}
        \partsitem Pemeriksaan aljabarnya mudah.
        \begin{equation*}
          0
          =xy_2+x_2y_3+x_3y-x_3y_2-xy_3-x_2y 
          =x\cdot (y_2-y_3)+y\cdot (x_3-x_2)+x_2y_3-x_3y_2 
        \end{equation*}
        dapat disederhanakan menjadi bentuk yang dikenal
        \begin{equation*}
          y=x\cdot (y_3-y_2)/(x_3-x_2)+(x_3y_2-x_2y_3)/(x_3-x_2)
        \end{equation*}
        (kasus $x_3-x_2=0$ mudah ditangani).

        Untuk memperoleh wawasan geometris, gambar ini
        menunjukkan kotak yang dibentuk oleh ketiga vektor tersebut.
        Perhatikan bahwa ujung ketiga vektor berada pada bidang $z=1$.
        Di bawah dua vektor di sebelah kanan terdapat garis yang melalui
        $(x_2,y_2)$ dan $(x_3,y_3)$.
        \begin{center}
          \includegraphics{det/mp/ch4.49}
        \end{center}
        Kotak itu akan memiliki volume tak nol kecuali jika segitiga yang
        dibentuk oleh ketiga ujung tersebut degenerat.
        Hal itu hanya terjadi (dengan mengandaikan
        $(x_2,y_2)\neq (x_3,y_3)$) jika $(x,y)$ terletak pada garis yang
        melalui kedua titik lainnya.
       \partsitem \answerasgiven %
        Kita mencari tinggi melalui $(x_1,y_1)$ dari segitiga dengan titik-titik
        sudut $(x_1,y_1)$, $(x_2,y_2)$, dan $(x_3,y_3)$ dengan cara biasa dari
        bentuk normal di atas:
        \begin{equation*}
          \frac{1}{\sqrt{(x_2-x_3)^2+(y_2-y_3)^2}}
          \begin{vmat}
            x_1  &x_2  &x_3  \\
            y_1  &y_2  &y_3  \\
            1    &1    &1
          \end{vmat}.
        \end{equation*}
        Langkah berikutnya menunjukkan bahwa luas segitiga adalah
        \begin{equation*}
          \frac{1}{2}\left|
          \begin{vmat}
            x_1  &x_2  &x_3  \\
            y_1  &y_2  &y_3  \\
            1    &1    &1
          \end{vmat}\right|.
        \end{equation*}
        Paparan ini mengungkapkan \textit{modus operandi} dengan lebih jelas
        daripada pembuktian biasa yang menunjukkan bahwa sekumpulan suku
        identik dengan determinan tersebut.
       \partsitem  \answerasgiven %
        Misalkan
        \begin{equation*}
          D=
          \begin{vmat}
            x_1  &x_2  &x_3  \\
            y_1  &y_2  &y_3  \\
            1    &1    &1
          \end{vmat}
        \end{equation*}
        maka luas segitiga adalah $(1/2)\deter{D}$.
        Jika semua koordinat berupa bilangan bulat, maka $D$ adalah bilangan bulat.
      \end{exparts}
    \end{answer}
\end{exercises}
